\begin{proof}
\begin{enumerate}[label={$\langle1\rangle$\arabic*.},leftmargin=*]
\item 
\noindent\textbf{Step 1.} Every irreducible component of a reduced hypersurface is generically smooth,
so $S_{\rm reg}$ is dense and its complement is a proper analytic subset.
If a local section of $\Omega^1_S$ vanishes on $S_{\rm reg}$, its support avoids
all minimal primes of the reduced local ring.  Prime avoidance therefore gives
an element of its annihilator outside every minimal prime; that element is a
nonzerodivisor, so the section is torsion.  Conversely,
$\Omega^1_S|_{S_{\rm reg}}$ is locally free, hence has no torsion, and every
torsion section restricts to zero there.

\end{enumerate}
\end{proof}

\begin{proof}
\begin{enumerate}[label={$\langle1\rangle$\arabic*.},leftmargin=*]
\item 
The conormal sequence is right exact.  It is also left exact here.  Locally
write $S=(g=0)$ with $g$ reduced.  If $h\,dg|_S=0$, then $h$ vanishes at the
generic smooth point of every irreducible component, because $dg$ is nonzero
there.  The ring $\mathcal O_S$ is reduced, so the intersection of its
minimal primes is zero; hence $h=0$.  Thus
\[
0\to N^*_{S/X}\to\Omega_X^1|_S\to\Omega_S^1\to0
\]
is exact.  By definition, $\mathscr K_S$ is the inverse image of
$\mathscr T_S$ under the last surjection.  Restriction to that inverse image
is onto $\mathscr T_S$ and has kernel $N^*_{S/X}$, proving the first
sequence.  Composing with the quotient by $\mathscr T_S$ proves the second.

\end{enumerate}
\end{proof}

\begin{proof}[The minimal-prime formulation of conormal injectivity]
\begin{enumerate}[label={$\langle1\rangle$\arabic*.},leftmargin=*]
\item 
Right exactness is the conormal sequence for a closed immersion.  We check
the left end.  Work at a point of $S$ and let $R$ be the regular local ring
of $X$, with reduced local equation $g$ for $S$.  The map at the left is
\[
 R/(g)\longrightarrow (\Omega^1_R/g\Omega^1_R),
 \qquad \bar h\longmapsto h\,dg.
 \tag{W.145}\hypertarget{tag-03DW}{}\label{wbII:eq:CN:2.2}
\]
Let $\mathfrak p$ be a minimal prime of $(g)$.  Since $R$ is regular and
has characteristic zero, the height-one local ring $R_{\mathfrak p}$ is a
discrete valuation ring.  Reducedness of $(g)$ gives
$\operatorname{ord}_{\mathfrak p}(g)=1$, so $g$ is a unit times a
uniformizer there; in characteristic zero the differential of that
uniformizer is nonzero modulo $\mathfrak p$.  Hence $dg$ is nonzero in
$\Omega^1_{R_{\mathfrak p}}\otimes k(\mathfrak p)$, so \eqref{wbII:eq:CN:2.2} sends $\bar h$
to zero only if $h\in\mathfrak p$.  This holds for every minimal prime.
Because $R/(g)$ is reduced, the intersection of its minimal primes is zero;
therefore $\bar h=0$.  Thus the conormal sequence in \cref{cdp:st:A.1.1} is exact.

\end{enumerate}
\end{proof}

\begin{proof}
\begin{enumerate}[label={$\langle1\rangle$\arabic*.},leftmargin=*]
\item 
Tensoring the second sequence of
\cref{cdp:st:A.1.1} by the locally free sheaf $A$
preserves exactness; its global-section sequence gives the first equality.
Tensor the first sequence and take global sections.  If the conormal
$H^0$ vanishes, exactness identifies the middle group with the displayed
kernel.  No condition on the torsion-free quotient controls that kernel.

\end{enumerate}
\end{proof}

\begin{proof}
\begin{enumerate}[label={$\langle1\rangle$\arabic*.},leftmargin=*]
\item 
The fibre ideal is the principal ideal $(g)$ and is reduced, so
$(g)/(g^2)\simeq\mathcal O_S$.  Multiplication by a unit does not affect
injectivity, and the term $f\,du$ vanishes after restriction to $f=0$.
It therefore suffices to use $f=x,xy,xyz$.

For $f=xy$, if a lift $\widetilde h$ satisfies
$\widetilde h(y\,dx+x\,dy)=0$ modulo $(xy)$, then
$\widetilde h y,\widetilde h x\in(xy)$.  Hence
$\widetilde h\in(x)\cap(y)=(xy)$ in the regular local ring, so $h=0$ on
$S$.  For $f=xyz$, the three coefficient conditions say
$\widetilde h yz,\widetilde h xz,\widetilde h xy\in(xyz)$; hence
$\widetilde h\in(x)\cap(y)\cap(z)=(xyz)$.  The case $f=x$ is immediate.
Thus the conormal map is injective in every local model, proving \eqref{wbII:eq:W:5.4}.

\end{enumerate}
\end{proof}

\begin{proof}
\begin{enumerate}[label={$\langle1\rangle$\arabic*.},leftmargin=*]
\item 
The assertion is local.  Away from $D$, $\eta$ is an isomorphism.
At a double point $S=\{xy=0\}$, a member of $\ker\beta$ is
\[
\bigl((h_1(y,z),h_2(x,z)),k(z)\bigr),
\qquad
h_1(0,z)=h_2(0,z)=k(z).
\]
The function
\[
h=h_1(y,z)+h_2(x,z)-h_1(0,z)
\]
on $\{xy=0\}$ restricts to $h_1,h_2$ on the two branches and to $k$ on
their intersection.  This proves middle exactness in the double model.

At a triple point $S=\{xyz=0\}$, write the branch functions as
$h_1(y,z)$, $h_2(x,z)$, $h_3(x,y)$.  Membership in $\ker\beta$ is exactly
\[
h_1(0,z)=h_2(0,z),\quad
h_1(y,0)=h_3(0,y),\quad
h_2(x,0)=h_3(x,0).
\]
Then
\[
\begin{split}
h={}&h_1(y,z)+h_2(x,z)+h_3(x,y)\\
&-h_1(0,z)-h_1(y,0)-h_2(x,0)+h_1(0,0)
\end{split}
\]
restricts to the prescribed function on every branch.  Injectivity of
$\alpha$ follows because $S$ is reduced and $\eta$ is an isomorphism over a
dense open set.  Surjectivity of $\beta$ follows locally by lifting a
function on the branch conductor to the smooth branch and taking $k=0$.
This proves exactness in every local model.

\end{enumerate}
\end{proof}
